\ExerciseSection

¶ 1) For every finite sequence \(s = (z_k)_{k \in \mathbf{I}}\) of complex numbers, put

\[
\sum_ {k \in \mathbf {I}} | z _ {k} | = \rho_ {s} \cdot \sup _ {\mathbf {J} \subset \mathbf {I}} \left| \sum_ {k \in \mathbf {J}} z _ {k} \right|
\]

(cf.~Chapter VII, § 3, no. 1, Proposition 2).

\begin{enumerate}
\def\labelenumi{\alph{enumi})}
\tightlist
\item
  * If \(z_{k} = r_{k} (\cos \varphi_{k} + i \sin \varphi_{k})\) , where \(r_{k} = |z_{k}|\) and \(\varphi_{k}\) is the principal radian measure of the amplitude of \(z_{k}\) , put
\end{enumerate}

\[
f (\varphi) = \sum_ {k \in \mathbf {I}} r _ {k} (\cos (\varphi - \varphi_ {k})) ^ {+}
\]

for every \(\varphi \in [0, 2\pi[\) . Show that the least upper bound of \(f(\varphi)\) in the interval \([0, 2\pi[\) is equal to \(\sup_{J \subset I} \left| \sum_{k \in J} z_k \right|\) . By considering the integral \(\int_0^{2\pi}f(\varphi)d\varphi\) , deduce that for all finite sequences \(s\) we have \(\rho_{s}\leqslant \pi\) and that there is no finite sequence \((z_k) = s\) such that \(\rho_{s} = \pi\)

\begin{enumerate}
\def\labelenumi{\alph{enumi})}
\setcounter{enumi}{1}
\tightlist
\item
  Show that, for every \(\varepsilon > 0\) , there exists a finite sequence \((z_k) = s\) such that \(\rho_s > \pi - \varepsilon\) (take the \(z_k\) to be the roots of a binomial equation of sufficiently high degree). *
\end{enumerate}

\begin{enumerate}
\def\labelenumi{\arabic{enumi})}
\setcounter{enumi}{1}
\tightlist
\item
  Let \((z_{n})\) be an infinite sequence of complex numbers \(z_{n} = x_{n} + iy_{n}\) such that \(x_{n} \geqslant 0\) for all \(n\) . Show that if the series whose general terms are \(z_{n}\) and \(z_{n}^{2}\) are convergent, then the series whose general term is \(z_{n}^{2}\) is absolutely convergent. Give an example in which this result is false when the condition \(x_{n} \geqslant 0\) (which can also be written \(-\pi / 2 \leqslant \operatorname{Am}(z_{n}) \leqslant \pi / 2\) ) is replaced by
\end{enumerate}

\[
- (\mathrm{I} - \varepsilon) \frac {\pi}{2} \leqslant \operatorname{Am} (z _ {n}) \leqslant (\mathrm{I} + \varepsilon) \frac {\pi}{2}
\]

no matter how small the real number \(\varepsilon > 0\) may be. (Here Am \((z)\) denotes that radian measure of the amplitude of \(z_{n}\) which belongs to the interval \([- \pi, +\pi]\) .)

\begin{enumerate}
\def\labelenumi{\arabic{enumi})}
\setcounter{enumi}{2}
\item
  Let \((z_{t})_{t\in I}\) be a family of complex numbers such that \(\prod_{t\in I}|z_t| = +\infty\) , (resp. \(\prod_{t\in I}|z_t| = 0\) ). Show that the family \((z_{t})\) is multipliable in the space \(\tilde{\mathbf{C}}\) ( \(\S 4\) , no. 3) and that its product is \(\infty\) (resp. 0).
\item
  Show that the infinite product whose general factor is \(1 + i/n\) is not convergent, but that the product of the absolute values of its factors is absolutely convergent.
\item
  Let \((z_{n})\) be an infinite sequence of non-zero complex numbers such that \(\lim_{n\to \infty}z_n = 1\) . Show that, if there exist permutations \(\sigma\) of \(\mathbf{N}\) such that the infinite product whose general factor is \(z_{\sigma (n)}\) is convergent, then the set of products
\end{enumerate}

\[
\underset {n = 0} {\overset {\infty} {\mathbf {P}}} z _ {\sigma (n)}
\]

corresponding to all these permutations can be (i) a single point; or (ii) the whole group \(\mathbf{C}^*\) ; or (iii) an open ray with origin \(\mathbf{o}\) ; or (iv) a circle with centre \(\mathbf{o}\) ; or (v) a ``logarithmic spiral'', the image of \(\mathbf{R}\) under the mapping \(t \to a^{t - t_0} (\cos t + i \sin t)\) , where \(a\) is a real number \(>0\) and \(\neq 1\) . (Argue as in Exercise 2 of Chapter VII, § 3, using the fact that the multiplicative group \(\mathbf{C}^*\) is isomorphic to the additive group \(\mathbf{R} \times \mathbf{T}\) .)

